% ============================================================
%  SECTION 3 — PROPOSED METHODOLOGY
%  For use in IEEE/ACM paper (paste into Overleaf body)
%  Required packages (add to preamble if not already present):
%    \usepackage{amsmath,amssymb}
%    \usepackage{booktabs}
%    \usepackage{algorithm}
%    \usepackage{algpseudocode}
%    \usepackage{graphicx}
%    \usepackage{listings}
% ============================================================

\section{Proposed Methodology}
\label{sec:methodology}

% ------------------------------------------------------------------
\subsection{System Architecture Overview}
\label{subsec:arch_overview}
% ------------------------------------------------------------------

The proposed framework is constructed upon three orthogonal tiers of
responsibility, a design invariant that governs every architectural
decision in the system.
\textbf{Tier~1} comprises the deterministic code layer responsible for
evidence collection (reconnaissance) and pipeline flow orchestration.
\textbf{Tier~2} encompasses the large language model (LLM) agents that
perform contextual reasoning and strategic planning within each bounded
phase.
\textbf{Tier~3} constitutes a second deterministic code layer---the
\emph{Proof Gate}---that produces the final exploitation verdict by
evaluating raw HTTP evidence, operating exclusively on structured data
rather than probabilistic model output.

This tripartite separation embodies the core epistemological principle
of the architecture: the ground truth of a security vulnerability is
established by the observed HTTP response corpus, not by the output of
any language model.
Formally, for an evidence set
$\mathcal{E} = \{e_1, e_2, \ldots, e_n\}$
where each $e_i \in \mathcal{H}$ denotes a recorded HTTP exchange, the
final verdict
$v \in \{\textsc{Exploited},\, \textsc{Failed},\, \textsc{InfoExposure}\}$
is computed by a deterministic function
$\Phi: \mathcal{H}^{n} \to \mathcal{V}$,
with no LLM inference participating in the mapping from
$\mathcal{E}$ to $v$.

The orchestration backbone is implemented as a compiled
\texttt{StateGraph} using LangGraph~\cite{langgraph2024}, wherein each
processing phase is encoded as a pure node function and every
control-flow transition is encoded as a deterministic conditional edge.
No LLM is granted authority over graph routing; the transition function
$\delta: S \times \Sigma \to S$
is evaluated exclusively on enumerated state values in the
\texttt{BugState} domain.

\paragraph{State machine formalization.}
The per-bug execution lifecycle is governed by a finite automaton
$\mathcal{M} = (Q, \Sigma, \delta, q_0, F)$ where:
\begin{itemize}
  \item $Q = \{\textsc{Queued},\; \textsc{Debating},\;
    \textsc{DebateApproved},\; \textsc{Executing},\;
    \textsc{Verifying},\; \textsc{Exploited},\;
    \textsc{NotExploited},\; \textsc{ProofQualityFail},\;
    \textsc{InfoExposureOnly},\; \textsc{Skipped},\;
    \textsc{Error}\}$
  \item $\Sigma = \{\texttt{START\_DEBATE},\; \texttt{APPROVED},\;
    \texttt{STOP},\; \texttt{INSUFFICIENT\_CONTEXT},\;
    \texttt{START\_EXEC},\; \texttt{EXEC\_DONE},\;
    \texttt{CONFIRMED},\; \texttt{PROOF\_FAIL},\;
    \texttt{RETRY\_DEBATE},\; \texttt{GIVE\_UP}\}$
  \item $\delta$ is the transition function defined by a declarative
    \texttt{TRANSITIONS} table; every legal $(q, \sigma) \to q'$ triple
    is explicitly enumerated, prohibiting any implicit state advancement
  \item $q_0 = \textsc{Queued}$ and
    $F = \{\textsc{Exploited},\; \textsc{NotExploited},\;
    \textsc{InfoExposureOnly},\; \textsc{Skipped},\; \textsc{Error}\}$
\end{itemize}

\paragraph{Frozen strategy invariant.}
Upon debate approval, four fields are frozen into an immutable
\texttt{StrategyDocument}
$\sigma = (\sigma_s,\, \sigma_g,\, \sigma_c,\, \sigma_q)$,
comprising the strategic approach~$\sigma_s$, step-by-step execution
guide~$\sigma_g$, success conditions~$\sigma_c$, and verification
questions~$\sigma_q$, respectively.
The execution agent receives $\sigma$ as read-only context and is
architecturally prohibited from modifying it.
Should the strategy prove deficient, the only corrective pathway is a
new debate round producing a revised $\sigma'$ after adversarial review.

% ------------------------------------------------------------------
\subsection{Pipeline and Data Flow Specifications}
\label{subsec:pipeline}
% ------------------------------------------------------------------

\subsubsection{Main Pipeline}

The top-level pipeline is compiled as a linear sequence of five nodes
executed unconditionally in order:
\begin{equation}
  \textsc{Start}
  \;\to\; r
  \;\to\; h
  \;\to\; k
  \;\to\; b
  \;\to\; \rho
  \;\to\; \textsc{End}
  \label{eq:main_pipeline}
\end{equation}
where $r$ denotes the reconnaissance node, $h$ the hypothesis generation
(Hunt) node, $k$ the coordination node, $b$ the per-bug execution loop,
and $\rho$ the report generation node.
Inter-node data flow is mediated by a shared
\texttt{PipelineState} TypedDict
$\mathcal{P} = (\mathcal{T},\, \mathcal{R},\, \mathcal{D},\, \mathcal{F})$
comprising the target specification $\mathcal{T}$, the reconnaissance
artifact $\mathcal{R}$, the ordered dossier list
$\mathcal{D} = [d_1, d_2, \ldots, d_m]$, and the confirmed finding set
$\mathcal{F}$.

\subsubsection{Per-Bug Sub-graph with Feedback Cycle}

For each dossier $d_i \in \mathcal{D}$, a dedicated sub-graph $G_i$ is
compiled and invoked.
The sub-graph admits a structured feedback cycle governed by the retry
predicate $\psi$:
\begin{align}
  G_i:\quad
  & \textsc{Debate}
    \xrightarrow{\textsc{Approved}}
    \textsc{Exec}
    \xrightarrow{\textsc{ExecDone}}
    \textsc{Verify}
  \notag \\
  & \textsc{Verify}
    \;\xrightarrow{\;\psi(\mathcal{E}_i)\;\wedge\;\textsc{ProofFail}\;}
    \textsc{Debate}
  \label{eq:bug_subgraph}
\end{align}
The retry predicate $\psi$ is evaluated as follows:
\begin{equation}
  \psi(\mathcal{E}_i)
  \;=\;
  \bigl(c_i \geq 1 \;\vee\; p_i\bigr)
  \;\wedge\;
  \bigl(v_i \leq v_{\max}\bigr)
  \;\wedge\;
  \bigl(r_i < 2\,r_{\max}\bigr)
  \label{eq:retry_predicate}
\end{equation}
where $c_i$ denotes the count of panel verifiers returning
\texttt{confirmed=True}; $p_i$ is a Boolean indicating partial proof
satisfaction (i.e., at least one \texttt{ProofMarker} with
\texttt{satisfied=True} in $\mathcal{E}_i$); $v_i$ and $v_{\max}$ are
the current and maximum verify-retry counts (default $v_{\max}=1$);
and $r_i$ and $r_{\max}$ are the current debate round count and the
configured maximum (default $r_{\max}=3$).
A per-bug wall-clock guard is enforced via
\texttt{asyncio.wait\_for} with a timeout of $\tau = 600$~s,
ensuring that no individual dossier can monopolize pipeline resources
regardless of LLM response latency.

\begin{figure}[htbp]
  \centering
  \includegraphics[width=\linewidth]{fig1_pipeline}
  % If the figure is not yet compiled, substitute the above with:
  % \fbox{\parbox{0.92\linewidth}{\centering\vspace{3.2cm}
  %   [Architecture Diagram Placeholder]\vspace{3.2cm}}}
  \caption{%
    End-to-end architecture of the proposed multi-agent web penetration
    testing framework.
    \textbf{Left (main pipeline):} The target URL and credentials enter
    the \textbf{Reconnaissance} node, which collects a
    \texttt{ReconArtifact}~$\mathcal{R}$ through four-layer HTTP
    crawling---passive link following, JavaScript endpoint extraction,
    active path probing against $\sim$50 curated generic paths plus
    dynamically inferred sibling IDOR candidates, and
    authenticated-versus-unauthenticated differential recording.
    The \textbf{Hunt} node invokes an LLM Hunter agent on $\mathcal{R}$,
    which emits an ordered list of \texttt{BugDossier} hypotheses
    $\mathcal{D} = [d_1,\ldots,d_m]$.
    The \textbf{Coordinate} node ranks $\mathcal{D}$ by severity,
    vulnerability category (BAC before BLF, simple before chain), and
    inter-dossier dependencies.
    The \textbf{Bugs} node iterates over $\mathcal{D}$, invoking the
    per-bug sub-graph for each $d_i$.
    The \textbf{Report} node synthesizes confirmed
    \texttt{Finding}~objects into \texttt{report.md},
    \texttt{findings.json}, and per-bug PoC scripts.
    \textbf{Right (per-bug sub-graph):} For each~$d_i$, the
    \textbf{Debate} phase runs an adversarial Red\,$\leftrightarrow$\,Blue
    agent loop, producing an immutable \texttt{StrategyDocument}~$\sigma$
    on \textsc{Approve}.
    The \textbf{Execution} phase instantiates a tool-calling LLM loop
    (up to 20~steps), recording every HTTP exchange via a
    \texttt{RecordingHttpClient} into an \texttt{Evidence}
    object~$\mathcal{E}_i$.
    The \textbf{Verification} phase runs three parallel
    \texttt{VerifierAgent} instances (advisory panel) followed by a
    deterministic \texttt{ProofGate}~$\Phi$.
    On \textsc{ProofQualityFail}, the gate's rejection rationale is fed
    back to the Debate node for a targeted strategy revision, subject to
    retry predicate~$\psi$ (Eq.~\ref{eq:retry_predicate}).
    All routing decisions are computed by deterministic code; no language
    model participates in control-flow decisions.
  }
  \label{fig:architecture}
\end{figure}

\subsubsection{Formal Data Definitions}

The primary data structures transferred between pipeline phases are
formally defined as follows.

The \textbf{ReconArtifact}
$\mathcal{R} = (E_{\mathrm{ep}},\, X_{\mathrm{http}},\, A_{\mathrm{prof}},\,
\Delta_{\mathrm{auth}},\, G_{\mathrm{wf}})$
encapsulates the discovered endpoint set~$E_{\mathrm{ep}}$, the observed
HTTP exchange corpus~$X_{\mathrm{http}}$, the collected authentication
profiles~$A_{\mathrm{prof}}$, the authentication-differential
set~$\Delta_{\mathrm{auth}}$ (pairs of anonymous versus authenticated
responses for each endpoint), and the inferred workflow
graph~$G_{\mathrm{wf}}$.

A \textbf{BugDossier}
$d = (\mathit{id},\, \kappa,\, \omega,\, \eta,\, \alpha,\, \mathrm{conf})$
specifies: the bug identifier; the vulnerability pattern code
$\kappa \in \{\textsc{BAC-01}, \textsc{BAC-02}, \textsc{BAC-03},
\textsc{BLF-01}, \textsc{BLF-05}, \textsc{BLF-06}, \ldots\}$;
the target endpoint~$\omega$; the attack hypothesis~$\eta$ in natural
language; the required actor configuration
$\alpha = (\mathit{role}_{\mathrm{attacker}},\,
\mathit{role}_{\mathrm{victim}},\,
\mathit{role}_{\mathrm{admin}})$;
and the confidence score $\mathrm{conf} \in [0,1]$.

An \textbf{Evidence} object
$\mathcal{E} = (X,\, \Pi,\, \Delta_s,\, C_{\mathrm{sess}})$
aggregates the recorded exchange sequence $X = [x_1, \ldots, x_n]$; the
proof-marker set $\Pi = \{\mu_1, \ldots, \mu_k\}$ where
$\mu_j = (\mathit{key}_j,\, \mathit{sat}_j,\, \mathit{detail}_j,\,
\mathit{seqs}_j)$;
the state delta~$\Delta_s$ mapping field names to numeric changes; and the
session-context map $C_{\mathrm{sess}}: \mathit{label} \mapsto \mathit{role}$.

Each \textbf{HttpExchange}
$x_i = (\mathrm{seq}_i,\, m_i,\, u_i,\, \hat{\omega}_i,\,
H_i^{\mathrm{req}},\, b_i^{\mathrm{req}},\, s_i,\,
H_i^{\mathrm{resp}},\, b_i^{\mathrm{resp}},\,
a_i,\, K_i,\, N_i,\, T_i)$
records the sequence number, HTTP method~$m_i$, URL~$u_i$, normalized
endpoint template~$\hat{\omega}_i$, request and response header/body
references, HTTP status code~$s_i$, actor label~$a_i$, JSON key set~$K_i$
(top 50 keys extracted per response), numeric field map~$N_i$, and HTML
title~$T_i$.
Body content is stored in a content-addressed \texttt{BodyStore}
(SHA-256-keyed binary files on disk); object references carry only a
hash pointer and a 200-character preview, ensuring that no response body
is truncated in transit between components.

% ------------------------------------------------------------------
\subsection{Core Components and Agent Execution}
\label{subsec:components}
% ------------------------------------------------------------------

\subsubsection{Reconnaissance Module}

The reconnaissance module operates as a four-layer evidence-collection
pipeline, each layer targeting a distinct category of endpoint discovery
not reachable by prior layers.

\textbf{Layer~1 (Passive crawl)} follows HTML hyperlinks and form
actions from within the authenticated session, capturing publicly linked
endpoints.
\textbf{Layer~2 (JavaScript extraction)} applies regular-expression
matching to inline \verb|<script>| blocks, recovering \texttt{fetch()},
\texttt{axios}, and \texttt{XMLHttpRequest} references not present in
the HTML DOM and reconstructing path templates from string-concatenation
patterns (e.g., \texttt{'/orders/' + id} $\to$ \texttt{/orders/\{id\}}).
\textbf{Layer~3 (Active probing)} issues HEAD/GET requests to
approximately 50 curated generic paths (administrative interfaces, REST
API prefixes, common commerce endpoints) augmented by dynamically
inferred IDOR candidate paths.
For each discovered collection endpoint~$\omega_c$ matching the pattern
\texttt{/api/v1/\{collection\}}, sibling resource probes
$\omega_c^{(j)} = \omega_c \cup \{{\omega_c}/{j} : j \in [1, N_p]\}$
are generated across 15 standardized REST collection names.
\textbf{Layer~4 (Authentication differential)} re-issues every
discovered request under both anonymous and authenticated sessions,
recording \texttt{AuthDiff} pairs
$(r_{\mathrm{anon}},\, r_{\mathrm{auth}})$ for each endpoint.
Endpoints satisfying
$\mathrm{status}(r_{\mathrm{anon}}) \in \{302, 403\}$
and
$\mathrm{status}(r_{\mathrm{auth}}) = 200$
constitute strong Broken Access Control signals.

A soft-404 filter prevents path inflation from applications that return
HTTP~200 with a ``page not found'' body.
For a probe response with body length~$\ell$, the probe is discarded if
$|\ell - \ell_{404}| < 100$ bytes, where $\ell_{404}$ is the body length
of a known-404 response.
JSON responses are unconditionally exempt from this filter: compact JSON
objects (e.g., 143 bytes) cannot be meaningfully compared to HTML error
pages (e.g., 2,800 bytes) via a byte-length heuristic.
An additional guard prevents requests to \texttt{/logout} (which would
terminate the authenticated session) and correctly interprets 302
redirects as ``endpoint exists, requires authentication'' rather than
``endpoint not found.''

Path normalization is applied via \texttt{\_family\_path}: numeric
segments and UUID-pattern segments in a URL path are replaced with the
canonical token \texttt{\{id\}}, enabling deduplication of structurally
identical endpoints discovered via different concrete resource
identifiers.

\subsubsection{Vulnerability Candidate Generation}

Candidate dossiers are produced through a two-stage pipeline.
A deterministic pre-filter applies rule-based nomination: endpoints
exhibiting non-empty authentication differentials are nominated as BAC
candidates; endpoints accepting numeric fields (\texttt{amount},
\texttt{quantity}, \texttt{price}) are nominated as BLF candidates;
endpoints with \texttt{\{id\}} path parameters are nominated as IDOR
(BAC-03) candidates.
An LLM Hunter agent then consumes the full \texttt{ReconArtifact}
alongside per-pattern signal digests provided by the
\texttt{PlaybookProvider} singleton---a module-level cached loader of
JSON vulnerability pattern cards and PortSwigger reference walkthroughs---
and emits a ranked list of \texttt{BugDossier} objects.

A post-processing reclassification pass enforces structural
correctness of pattern codes: collection-level endpoints (e.g.,
\texttt{/api/v1/users}) are promoted from BAC-03 to BAC-01;
per-object endpoints on non-administrative paths (e.g.,
\texttt{/api/v1/orders/\{id\}}) are demoted from BAC-01 to BAC-03.
Final deduplication is performed by the composite key
$(\hat{\omega},\, \kappa)$, where $\hat{\omega}$ denotes the normalized
endpoint family.

\subsubsection{Adversarial Debate Mechanism}

The debate phase instantiates an adversarial planning protocol between
two specialized agents: a \textbf{Red agent} (offensive strategist,
temperature $T_R = 0.6$) and a \textbf{Blue agent} (skeptical reviewer,
temperature $T_B = 0.3$).
The protocol runs for at most $r_{\max}$ rounds, as detailed in
Algorithm~\ref{alg:debate}.

\begin{algorithm}[htbp]
\caption{Adversarial Debate Protocol}
\label{alg:debate}
\begin{algorithmic}[1]
\Require Dossier $d$, \texttt{ReconArtifact} $\mathcal{R}$,
         memory context $\mathcal{M}_{\mathrm{ctx}}$,
         $r_{\max} = 3$
\Ensure Frozen strategy $\sigma$ or termination verdict
\For{$r \gets 0$ \textbf{to} $r_{\max} - 1$}
  \State $\sigma^{(r)} \gets \textsc{RedAgent.Argue}(d,\, \mathcal{R},\,
         \mathcal{M}_{\mathrm{ctx}},\, r)$
  \If{``\texttt{INSUFFICIENT\_EVIDENCE}'' $\in$ prefix$(\sigma^{(r)}, 500)$}
    \State \Return \textsc{InsufficientContext}
  \EndIf
  \State $v^{(r)} \gets \textsc{BlueAgent.Review}(\sigma^{(r)},\, \mathcal{R},\, r)$
  \If{$r = 0$ \textbf{and} $v^{(r)} = \textsc{Approve}$}
    \State $v^{(r)} \gets \textsc{Revise}$ \Comment{Mandatory conservatism constraint}
  \EndIf
  \If{\textsc{FieldGrounding}$(\sigma^{(r)}, \mathcal{R})$ fails}
    \State $v^{(r)} \gets \textsc{Revise}$ \Comment{Deterministic field-name objection}
  \EndIf
  \If{$v^{(r)} = \textsc{Approve}$}
    \State Freeze $(\sigma_s, \sigma_g, \sigma_c, \sigma_q) \gets \textsc{Parse}(\sigma^{(r)})$
    \State \Return \textsc{Approved}, $\sigma$
  \ElsIf{$v^{(r)} \in \{\textsc{Stop}, \textsc{Unverifiable}\}$}
    \State \Return \textsc{InsufficientContext}
  \EndIf
\EndFor
\State \Return \textsc{MaxRoundsExceeded}
\end{algorithmic}
\end{algorithm}

In round $r = 0$, the Red agent synthesizes an initial strategy by
reasoning over the dossier~$d$, the pattern playbook card for~$\kappa$,
the PortSwigger reference solution for~$\kappa$, the observed
authentication cookies from~$\mathcal{R}$, and any applicable
long-term memory techniques.
In rounds $r \geq 1$, the Red agent operates in rebuttal mode,
responding to exactly the objections raised by Blue in round $r-1$
rather than regenerating the strategy from scratch.

The mandatory conservatism constraint at line~8 ensures that at least
one adversarial exchange occurs before execution proceeds.
The deterministic field-grounding check at line~10 extracts the
money/quantity field names claimed in $\sigma^{(r)}$ from
backtick-delimited or quoted tokens and compares them against the known
form fields observed in~$\mathcal{R}$; a mismatch triggers an automatic
\textsc{Revise} override with a specific field-name objection,
independent of the Blue agent's LLM verdict.

\subsubsection{Execution Engine}

The execution engine implements a bounded tool-calling agent loop.
A \texttt{RecordingHttpClient} instance---a thin wrapper over
\texttt{httpx.AsyncClient}---intercepts all HTTP traffic issued by the
LLM agent, recording each request-response pair as a structured
\texttt{HttpExchange}~$x_i$ with full body preservation in the
\texttt{BodyStore}.
The \texttt{follow\_redirects} flag is set to \texttt{False}, ensuring
that 302 redirect responses are recorded as distinct evidence rather
than silently absorbed into the subsequent 200 response.

The agent is allocated a step budget of
$N_{\max} \in \{12,\, 20\}$ steps, where $N_{\max} = 20$ applies to
chained multi-step BLF patterns and $N_{\max} = 12$ applies to
simpler BAC patterns.
A sliding context window of depth $w = 8$ most-recent turns is
maintained to prevent token exhaustion while preserving local
action-response continuity; the system prompt and the first user turn
are always retained outside the window.

The execution prompt is augmented with four harvested context
categories:
(1)~\textbf{Endpoint schema}---form field names and data types inferred
from HTML forms, JSON POST response bodies, GET response schemas, and
cross-page forms;
(2)~\textbf{Known values}---concrete resource identifiers, coupon codes,
usernames, and product prices extracted from~$\mathcal{R}$, labeled
per-endpoint to prevent cross-endpoint value injection;
(3)~\textbf{Chain-steering directives}---for multi-step BLF patterns,
the \texttt{exploit\_approach} of~$d$ is injected with an explicit
step-ordering constraint;
(4)~\textbf{Memory skills}---abstract cross-target techniques promoted
from the long-term memory store.

The \texttt{RecordingHttpClient} implements automatic encoding
negotiation: if an HTTP POST/PUT/PATCH returns a status code in
$\{400, 404, 415, 422\}$, the request body is automatically
re-encoded---switching between \texttt{application/x-www-form-urlencoded}
and \texttt{application/json}---and retried, shielding the execution
agent from content-type negotiation failures.
An anti-stall guard terminates the loop if $s \geq 7$ steps have elapsed
with zero state-changing requests (POST/PUT/PATCH/DELETE returning 2xx),
preventing unbounded GET-only exploration.

When the LLM tool-calling loop produces no exploitation evidence,
two deterministic fallback routines are invoked.
\texttt{\_deterministic\_bac02\_attempt()} applies all viable
cookie-tampering combinations identified in~$\mathcal{R}$, flipping
\texttt{role}, \texttt{is\_admin}, and \texttt{isadmin} field values
(\texttt{false}$\to$\texttt{true}, \texttt{0}$\to$\texttt{1}).
\texttt{\_deterministic\_blf\_attempt()} constructs the maximal tampered
POST body from the harvested endpoint schema, setting every money or
quantity field to a negative value ($-100$ or $-1$, respectively).

\subsubsection{Evidence-Driven Verification Subsystem}

Verification is executed in two sequential stages: an advisory LLM
panel followed by a deterministic proof gate.
The complete procedure is summarized in
Algorithm~\ref{alg:verification}.

\begin{algorithm}[htbp]
\caption{Two-Stage Verification}
\label{alg:verification}
\begin{algorithmic}[1]
\Require Evidence $\mathcal{E}$, pattern $\kappa$,
         frozen questions $\sigma_q$, panel size $n = 3$
\Ensure Bug status $s^* \in \mathcal{V}$
\State $[V_1, V_2, V_3] \gets \texttt{asyncio.gather}(
       \textsc{Verify}(\mathcal{E}, \sigma_q, T{=}0.0),
       \textsc{Verify}(\mathcal{E}, \sigma_q, T{=}0.1),
       \textsc{Verify}(\mathcal{E}, \sigma_q, T{=}0.2))$
       \Comment{Parallel dispatch}
\State $c \gets |\{j : V_j.\mathit{confirmed} = \text{True}\}|$
\State $Q^* \gets \{q_i : \textstyle\sum_j V_j.Q[i] > n/2\}$
\State $\lambda_{\mathrm{llm}} \gets (|\sigma_q| \geq 2)
       \wedge (|Q^*| \geq 2)$
\State $\mathit{has2xx} \gets \exists\, x_i \in \mathcal{E} :
       200 \leq s_i < 300$
\If{$c = 0$ \textbf{and} \textbf{not} $\mathit{has2xx}$}
  \State \Return \textsc{ProofQualityFail}
  \Comment{Fast-fail: no usable evidence}
\EndIf
\State $v \gets \Phi(\mathcal{E},\, \kappa)$
       \Comment{Deterministic ProofGate evaluation}
\If{$v = \textsc{Exploited}$}
  \State \Return \textsc{Exploited}
\ElsIf{$\lambda_{\mathrm{llm}}$}
  \State \Return \textsc{Exploited}
  \Comment{OR path: panel passes when gate is inconclusive}
\ElsIf{$v = \textsc{InfoExposureOnly}$}
  \State \Return \textsc{InfoExposureOnly}
\Else
  \State \Return \textsc{ProofQualityFail}
\EndIf
\end{algorithmic}
\end{algorithm}

Three \texttt{VerifierAgent} instances are dispatched in parallel via
\texttt{asyncio.gather}.
Each agent is assigned a distinct epistemic lens---\textit{technical}
($T = 0.0$), \textit{impact} ($T = 0.1$), and \textit{skeptic}
($T = 0.2$)---and is intentionally isolated from the debate transcript
to prevent anchoring bias.
Each agent independently reviews up to 20 HTTP exchanges with response
bodies of up to 20,000 characters each and returns a
\texttt{VerifierVerdict} tuple
$(b_j,\, \gamma_j,\, \rho_j,\, M_j,\, Q_j)$
comprising: confirmation Boolean~$b_j$, confidence
$\gamma_j \in [0,1]$, rationale~$\rho_j$, cited proof-marker
keys~$M_j$, and per-question Boolean answers~$Q_j$.
The LLM question-pass predicate is:
\begin{equation}
  \lambda_{\mathrm{llm}}
  = \bigl(|\sigma_q| \geq 2\bigr)
  \;\wedge\;
  \bigl(|Q^*| \geq 2\bigr),
  \qquad
  Q^* = \Bigl\{q_i : \textstyle\sum_j Q_j[i] > \tfrac{n}{2}\Bigr\}
  \label{eq:llm_pass}
\end{equation}

The fast-fail condition at lines~7--8 of Algorithm~\ref{alg:verification}
bypasses the \texttt{ProofGate} when zero panel confirmations coincide
with the complete absence of HTTP 2xx-status exchanges in~$\mathcal{E}$,
indicating that no server-side action was successfully completed.
Otherwise, the deterministic \texttt{ProofGate}~$\Phi$ is evaluated
regardless of the panel count.
This design ensures that ProofGate serves as a safety net even when
the panel is unanimously skeptical but physical HTTP evidence exists.

\subsubsection{Proof Gate Specifications}

The \texttt{ProofGate} class hierarchy implements pattern-specific
exploitation criteria as deterministic membership functions over the
evidence corpus~$\mathcal{E}$.
Table~\ref{tab:proof_rules} summarizes the required
\texttt{ProofKey} markers and acceptance conditions for each supported
vulnerability pattern.

\begin{table}[htbp]
\centering
\caption{ProofGate evaluation criteria per vulnerability pattern.
All required markers must be simultaneously satisfied for an
\textsc{Exploited} verdict; partial marker satisfaction yields
\textsc{InfoExposureOnly}; no marker satisfaction yields
\textsc{Failed}.}
\label{tab:proof_rules}
\begin{tabular}{p{1.6cm} p{3.0cm} p{3.8cm}}
\toprule
\textbf{Pattern} & \textbf{Required Markers} & \textbf{Acceptance Condition} \\
\midrule
BAC-01 & \textsc{PrivilegedAccess} &
  Anonymous actor receives HTTP 200 with $\geq 2$ records
  containing sensitive fields (email, balance, role, \ldots);
  or LLM classifier confirms \texttt{exposes\_other\_users\_pii=True} \\[4pt]
BAC-02 & \textsc{PrivilegedAccess}, \textsc{AuthBypass} &
  Admin-titled page rendered for unprivileged actor;
  order-independent detection of cookie difference on the same
  endpoint (blocked $\to$ 200 after cookie tamper) \\[4pt]
BAC-03 & \textsc{OwnershipBypass} &
  Response contains owner field $\neq$ attacker user-ID;
  or $\geq 2$ exchanges on the same endpoint template yield
  distinct identity values (\texttt{\_HTML\_ID\_PATTERNS} fallback) \\[4pt]
BLF-01 & \textsc{PriceManipulation}, \textsc{StateDelta} &
  Negative/zero numeric value in money field accepted
  (2xx/3xx without error-marker tokens in body) \\[4pt]
BLF-05 & \textsc{PriceManipulation}, \textsc{StateDelta} &
  Same coupon code accepted $\geq 2$ times with a consume event
  (checkout/purchase/pay) occurring between the two uses;
  or $\geq 2$ distinct codes stacked; or
  grand-total $\leq 0$ in response \\[4pt]
BLF-06 & \textsc{QuantityTamper}, \textsc{StateDelta} &
  Negative quantity accepted with 2xx/3xx lacking error-marker tokens \\[4pt]
BLF-03 & \textsc{StateSkip}, \textsc{StateDelta} &
  $\geq 2$ distinct action endpoints (POST/PUT/PATCH/DELETE)
  return 2xx with a non-empty state delta \\
\bottomrule
\end{tabular}
\end{table}

An LLM-assisted classifier (\texttt{ClassifyResponses}) is invoked as a
feature-extraction fallback when direct field-matching fails.
The classifier receives up to three representative 2xx exchanges,
selected by descending response body length and deduplicated by the
pair $(u_i, a_i)$ (URL and actor label), bounded at
$\ell_{\max} = 6{,}000$ characters.
Per-pattern JSON schemata constrain the classifier output to structured
fact objects (e.g., for BAC-03: \texttt{\{owner\_field\_name,
owner\_field\_value, email\_in\_response\}}).
Temperature is set to $T = 0.0$ on the initial attempt and escalated
to $T = 0.4$ on retries (up to three total attempts).
The extracted facts are evaluated by the same deterministic gate rules,
ensuring that the LLM functions solely as a feature extractor rather
than as a verdict authority.

\subsubsection{Long-Term Memory Subsystem}

A two-tier memory architecture is maintained in a local SQLite store to
accumulate and generalize exploitation knowledge across penetration
testing sessions.

\textbf{Tier~1 (Episodic memory)} records each exploitation attempt as
an \texttt{Episode} row keyed by a target fingerprint
$\phi = \mathrm{hash}(\mathit{host} \;\|\; \mathit{title} \;\|\;
\mathrm{sorted}(E_{\mathrm{ep}}))$.
The payload~$p$, proof-marker set~$\Pi$, and outcome are persisted for
every completed dossier execution.
A write gate enforces that only episodes in which the
\texttt{ProofGate} returned \textsc{Exploited} are written to the
store; advisory LLM panel opinions do not trigger writes,
preventing unverified hypotheses from polluting the memory corpus.

\textbf{Tier~2 (Semantic memory)} promotes abstract techniques to a
\texttt{rules} table when the accumulated episode evidence satisfies the
promotion predicate:
\begin{equation}
  \mathit{promote}(\kappa, f)
  = \bigl(|\mathcal{S}_{\kappa,f}| \geq \theta_s\bigr)
  \;\wedge\;
  \bigl(\bigl|\{\phi : (\kappa, f, \phi) \in \mathcal{S}_{\kappa,f}\}\bigr|
  \geq \theta_t\bigr)
  \label{eq:ltm_promote}
\end{equation}
where $\mathcal{S}_{\kappa,f}$ is the set of successful episodes for
pattern~$\kappa$ and endpoint family~$f$, $\theta_s = 3$ is the minimum
success count, and $\theta_t = 2$ is the minimum distinct-target count.

Promotion is withheld until a technique has been independently confirmed
on at least $\theta_t = 2$ distinct application fingerprints, preventing
single-application-specific payloads from being treated as general
techniques.
Endpoint paths are normalized to a \emph{role label}
(e.g., \texttt{/wallet/transfer} and \texttt{/account/send} both resolve
to the \texttt{transfer} role) before promotion, enabling cross-target
generalization across applications with structurally equivalent
endpoints but different path names.
Promoted rules are serialized as abstract, host-stripped descriptions
and injected into subsequent execution prompts as supplementary
\textit{Past Successful Techniques}.
